\section{Related Work}
\label{sec:related}

This section positions the work against its neighbours and states the delta at the width the evidence
supports.

\subsection{Measurement bias, and the harness as its current object}

The shape of our finding is old in systems research and new only in the evaluation of agents: that a
measurement depends on the stack under it is unremarkable in performance work, where declaring the
configuration is a norm. What agentic evaluation lacks is that norm. The argument has been made for the
harness as a whole~\citep{zhang2026harness,yao2026harnessbench}; this study is downstream of it and supplies the
measurement it lacks --- what the undeclared choices are worth, one axis at a time.

Four lines of work bound this one, and the artifact expands each with the figures we take from it.
\textbf{Apparatus variance is established methodology elsewhere}: rigorous performance measurement has
long required reporting the configuration and treating run-to-run variation as first
class~\citep{georges2007rigorous,iosup2011variability}, and measurement bias from unreported setup
choices was demonstrated two decades ago~\citep{mytkowicz2009wrongdata} --- what is missing here is not
the method but its application. \textbf{Non-determinism at nominal temperature zero} has been measured for repeated
sampling~\citep{atil-nondeterminism} and traced to the batch size a server happens to be
running~\citep{he2025nondeterminism} --- the class of server-side quantity our \S\ref{sec:runtime} finds
differing between two services at fixed model identity. \textbf{The endpoint is already known to be a service rather than an artifact}: a measurement study of
hosted open-weight APIs finds one model differing across providers in protocol support, error semantics,
context capacity and throughput~\citep{li2026service}, which is our T5--T8 axis arrived at independently.
\textbf{And the transport axis is visible to the field, in no fixed direction}: a leaderboard ranks every
model in a native and a prompted mode~\citep{bfcl2025}; on one model the prompted variant is \emph{ahead}
by 6pp~\citep{bfcl-leaderboard,bfcl-issue606}; on an agentic tool-use benchmark a prompt-based protocol
severely degrades a model of the same family as our T3/T4, attributed to a template mismatch with the
model's training~\citep{mcpverse2025}; and a third transport, tools as typed code stubs rather than JSON,
matches or exceeds native calling on 11 of 14 models~\citep{bitterlesson2026}. Three transports, and no
direction that survives a change of task family --- which is what makes the axis a parameter to declare
rather than a correction to apply. \textbf{And the task family has its own literature}:
reconstruction from stripped binaries is active~\citep{tan2024llm4decompile,jiang2025nova}, semantic
equivalence under LLM-driven transformation has been evaluated directly, on a different corpus and
task~\citep{detomasi2025obfuscation}, test-based equivalence has known
limits~\citep{liu-wang-issta2020}, and the reproducibility of published results built on
commercial models has been assessed directly, with model deprecation as a leading cause of
failure~\citep{angermeir2026reproducibility}.

\subsection{The transport: known to move the number, in no fixed direction}

The nearest thing to a counterexample is the leaderboard above: it ranks every model in both a native and a
prompted mode, so the axis is not invisible to the field. It reports the two modes side by side rather than
crossing them with the serving endpoint, and on one model the prompted variant is \emph{ahead} by
6pp~\citep{bfcl-leaderboard,bfcl-issue606} --- which we treat as a constraint on our claim rather than as
support. Designing on an effect size estimated from earlier data is known to miss its target
power~\citep{albers2018power}; what this case adds is a direction that reverses rather than drifts.

The reporting obligation is already established and we do not argue for it: community guidelines for
empirical work with language models require the harness to be described when it goes beyond bare API
calls~\citep{baltes2026guidelines}. Those guidelines do not make all of the apparatus fields examined here
equally explicit, and Table~\ref{tab:checklist} says which. What this study supplies is the other half: how much the
choices in question move a measured number.

\subsection{The endpoint: a unit of measurement, and a unit of selection}

\textbf{That two services can serve one model identity and return different output distributions is
established, and measured at scale.} A two-sample test on the Maximum Mean Discrepancy, applied to
commercial inference APIs for four Llama models, finds \textbf{11 of 31 endpoints serving distributions
that differ from the reference weights} --- naming Amazon Bedrock and Azure AI Studio among the access
paths, at a median 77.4\% power against a range of distortions~\citep{gao2025equality}. That work settles
the premise this study starts from, and it does so on raw output distributions from single prompts: it
crosses no transport, runs no multi-turn agentic task, and says nothing about whether the divergence
reaches a downstream score. What \S\ref{sec:results} adds is one case where it reaches the
\emph{stability} of such a score.

Holding weights, decoding and hardware fixed while varying only the inference engine produces measurable
differences in output, which establishes the mechanism this study extends: if an engine moves a number at fixed
weights, a managed platform choosing engine, version and scaffolding without disclosing any of them can
move it too, and a reader has no way to know by how much~\citep{silent-hyperparameter}.

Removal also breaks reproduction downstream, and its scale has been measured on one public arena's
own roster: of 243 models, 205 had been silently deprecated against 47 marked as such, at rates of
87.8\% for open-weight and 80\% for proprietary models~\citep{leaderboard-illusion}. That audit
names the mechanism --- ``Chatbot Arena may be forced to deprecate a model when a provider no longer
supports it via its API'' --- and counts the outcome after the fact, aggregated over a roster; this census measures the refusals themselves, verbatim and case by case, before a score exists.

\subsection{What the census adds to the failure taxonomies}

The nearest neighbour to our census is a taxonomy of failure modes in multi-provider serving
infrastructure, classified by origin layer and detectability~\citep{failureatlas}. Two of its
five layers already reach failures occurring before generation --- network and transport covers
``TCP timeouts, TLS handshake failures, DNS resolution issues'', governance and cost covers
``rate-limit counters that leak and never recover'' --- so the distinction we draw is not about
\emph{when} the failure lands.

What the census adds is a different question, asked one step earlier: not how a request fails once a
model is being served, but \emph{whether this endpoint will serve this model identity at all}, under this
request contract, from this account and this region. Several of our mechanisms are access and lifecycle
decisions rather than defects, and those persist for the duration of a study whoever is judged
responsible; others --- the request-contract constraints and the adapter incompatibility --- could
reasonably be filed as bugs by an operator who owns them. \textbf{Nothing in our contribution depends on
which filing a reader prefers}: what \S\ref{sec:tassonomia} establishes is that mechanisms of this class
exist, where in the evaluation lifecycle they act, what they cost a measurement, and that they need
different detection. The artifact carries the layer-by-layer comparison for a reader who wants to place
each row.

